\documentclass[12pt,a4paper]{article}
\usepackage[spanish]{babel}
\usepackage[utf8]{inputenc}
\usepackage[T1]{fontenc}
\usepackage{lmodern}
\usepackage[table]{xcolor}
\usepackage{geometry}
\usepackage{setspace}
\usepackage{longtable}
\usepackage{booktabs}
\usepackage{array}
\usepackage{tabularx}
\usepackage{hyperref}

\geometry{margin=2.4cm}
\setstretch{1.12}
\usepackage{enumitem}


\definecolor{PrimaryBlue}{HTML}{16324F}
\definecolor{SoftGray}{HTML}{F3F6F8}

\title{\Huge\textbf{Sprint Review}\\[0.3em]\Large Plataforma Web de Mascotas 3D}
\author{Bryan Quispe}
\date{20 de abril de 2026}

\begin{document}

\begin{titlepage}
\centering
\vspace*{1.8cm}
{\Huge\bfseries\color{PrimaryBlue} Sprint Review\\[0.35em]}
{\LARGE Revisión del incremento y retroalimentación\\[1.2cm]}

\begin{tabular}{p{5cm} p{8cm}}
\rowcolor{SoftGray}
\textbf{Versión} & 1.0 \\
\textbf{Fecha} & 20 de abril de 2026 \\
\textbf{Autor} & Bryan Quispe \\
\textbf{Periodo} & 20 de abril de 2026 al 7 de agosto de 2026 \\
\end{tabular}

\vfill
{\large Documento Scrum para la revisión del incremento}
\end{titlepage}

\section{Incremento mostrado}
\begin{itemize}[leftmargin=*]
    \item Registro de usuario.
    \item Inicio de sesión con JWT.
    \item Alta de mascotas.
    \item Lista de mascotas por usuario.
\end{itemize}

\section{Feedback}
\begin{itemize}[leftmargin=*]
    \item Mejorar mensajes de error para usuario final.
    \item Aclarar la expiración por inactividad.
    \item Mantener consistencia visual en formularios.
\end{itemize}

\section{Decisiones}
\begin{itemize}[leftmargin=*]
    \item Aceptar el incremento base del Sprint 1.
    \item Registrar mejoras visuales como nuevas tareas del siguiente Sprint.
\end{itemize}

\end{document}
